\begin{capa}
{\linespread{1.15}\selectfont
\begin{center}
INSTITUTO FEDERAL DE EDUCAÇÃO, CIÊNCIA E TECNOLOGIA DE GOIÁS -- CÂMPUS URUAÇU\par
DEPARTAMENTO DE ÁREAS ACADÊMICAS\par
CURSO DE TECNOLOGIA EM ANÁLISE E DESENVOLVIMENTO DE SISTEMAS\par
\vspace*{110pt}
EDUARDO HENRIQUE GUIMARÃES MARQUES\par
\vspace*{130pt}
\textbf{DESENVOLVIMENTO DE PLATAFORMA WEB PARA APOIO À ANÁLISE DE EEG COM INTELIGÊNCIA ARTIFICIAL EXPLICÁVEL}\par
\vfill
URUAÇU/GO\par
2026\par
\end{center}
}
\end{capa}

\begin{folhaderosto}
{\linespread{1.15}\selectfont
\begin{center}
EDUARDO HENRIQUE GUIMARÃES MARQUES\par
\vspace*{120pt}
\textbf{DESENVOLVIMENTO DE PLATAFORMA WEB PARA APOIO À ANÁLISE DE EEG COM INTELIGÊNCIA ARTIFICIAL EXPLICÁVEL}\par
\vspace*{90pt}
\hspace*{8.5cm}\begin{minipage}{7.2cm}
\fontsize{10}{12}\selectfont
\begin{SingleSpace}
Trabalho de Conclusão de Curso apresentado ao curso de Tecnologia em Análise e Desenvolvimento de Sistemas do Departamento de Áreas Acadêmicas do Instituto Federal de Educação, Ciência e Tecnologia de Goiás -- Câmpus Uruaçu como requisito parcial para conclusão do curso.

\vspace{0.8cm}
Orientador: Dr. Lynwood Livi de Souza
\end{SingleSpace}
\end{minipage}
\par\vfill
URUAÇU/GO\par
2026\par
\end{center}
}
\end{folhaderosto}

\begin{agradecimentos}
Agradeço primeiramente a Deus, pela saúde, força e perseverança concedidas ao longo de toda a minha caminhada acadêmica.

Aos meus pais e à minha família, pelo apoio incondicional, compreensão e incentivo constante em cada etapa da minha formação.

Ao meu orientador, Prof. Dr. Lynwood Livi de Souza, pela orientação criteriosa, paciência, disponibilidade e pelas valiosas contribuições técnicas e científicas que guiaram a evolução deste trabalho desde a iniciação tecnológica até a conclusão deste curso.

Ao Instituto Federal de Educação, Ciência e Tecnologia de Goiás (IFG) – Câmpus Uruaçu, e a todos os professores do Curso de Tecnologia em Análise e Desenvolvimento de Sistemas, pelos ensinamentos transmitidos e pela dedicação à educação pública e de qualidade.

Aos colegas de curso e amigos que compartilharam desafios, aprendizados e momentos importantes durante a graduação, e aos membros da banca examinadora, pela disponibilidade e pelas contribuições para o aperfeiçoamento deste trabalho.
\end{agradecimentos}

\begin{resumo}
A interpretação de exames de Eletroencefalograma (EEG) no diagnóstico da epilepsia exige a inspeção visual de longas séries temporais multicanais, processo sujeito à fadiga e à variabilidade entre avaliadores. Embora modelos de aprendizado profundo possam reconhecer padrões complexos, sua aplicação clínica enfrenta a fragmentação entre ferramentas e a dificuldade de inspecionar as decisões das redes neurais. Este trabalho desenvolveu e avaliou o protótipo funcional de uma plataforma web integrada (EEG-XAI) voltada ao apoio à triagem e à análise neurofisiológica. A solução adota arquitetura cliente-servidor mediada por API REST e combina uma interface em React, TypeScript e Recharts com um núcleo de processamento em Python, FastAPI, MNE-Python, TensorFlow/Keras e PostgreSQL. No motor preditivo, implementou-se a arquitetura híbrida CNN-BiLSTM v43, com fusão tardia de 23 canais e sequências de 15 características tempo-frequência por canal, incluindo entropia espectral. Para aumentar a inspecionabilidade, o Gradient SHAP foi integrado ao pipeline e suas relevâncias positivas foram projetadas sobre o traçado, por canal e intervalo temporal, para a sequência de pico. Nos experimentos com registros contínuos da base pública CHB-MIT, os grupos de treino e calibração foram separados antes do janelamento. Na calibração, o modelo atingiu F1 por EDF de 95,7\%, F1 localizado de 90,9\% e detecção de 15 em 20 eventos ictais, com 0,55 falso alarme por hora. Na avaliação histórica de desenvolvimento, obteve F1 por EDF de 82,4\%, F1 localizado de 46,2\% e detecção de 4 em 9 eventos, evidenciando a dificuldade de generalização interpaciente. Conclui-se que a plataforma integra ingestão de EDF, navegação multicanal, localização temporal de trechos suspeitos, explicabilidade visual e emissão de laudos, oferecendo apoio computacional à revisão médica sem substituir o diagnóstico clínico.

\vspace{\onelineskip}
\noindent\textbf{Palavras-chave}: eletroencefalograma; epilepsia; aprendizado profundo; CNN-BiLSTM; inteligência artificial explicável; SHAP.
\end{resumo}

\begin{resumo}[ABSTRACT]
\begin{otherlanguage*}{english}
Interpreting Electroencephalogram (EEG) exams for epilepsy diagnosis requires visual inspection of long multichannel time series, a process subject to fatigue and inter-observer variability. Although deep learning models can recognize complex patterns, their clinical application faces fragmented workflows and limited decision inspectability. This work developed and evaluated a functional integrated web platform prototype (EEG-XAI) to support neurophysiological triage and analysis. The solution adopts a REST API-mediated client-server architecture and combines a React, TypeScript, and Recharts interface with a processing core based on Python, FastAPI, MNE-Python, TensorFlow/Keras, and PostgreSQL. Its predictive engine implements the v43 hybrid CNN-BiLSTM architecture with late fusion across 23 channels and sequences of 15 time-frequency features per channel, including spectral entropy. To increase inspectability, Gradient SHAP was integrated into the pipeline, and positive relevance values for the peak sequence were projected onto the trace by channel and time interval. In experiments with continuous recordings from the public CHB-MIT dataset, training and calibration groups were separated before windowing. In calibration, the model achieved an EDF-level F1-score of 95.7\%, a localized F1-score of 90.9\%, and detected 15 of 20 ictal events, with 0.55 false alarms per hour. In the historical development evaluation, it achieved an EDF-level F1-score of 82.4\%, a localized F1-score of 46.2\%, and detected 4 of 9 events, highlighting the difficulty of inter-patient generalization. The platform integrates EDF ingestion, multichannel navigation, temporal localization of suspicious segments, visual explainability, and medical reporting, providing computational support for specialist review without replacing clinical diagnosis.

\vspace{\onelineskip}
\noindent\textbf{Keywords}: electroencephalogram; epilepsy; deep learning; CNN-BiLSTM; explainable artificial intelligence; SHAP.
\end{otherlanguage*}
\end{resumo}

\pdfbookmark[0]{\listfigurename}{lof}
\listoffigures*
\cleardoublepage
\pdfbookmark[0]{LISTA DE QUADROS}{loq}
\listoftables*
\cleardoublepage

\begin{siglas}
\item[ANVISA] Agência Nacional de Vigilância Sanitária
\item[API] Application Programming Interface (Interface de Programação de Aplicações)
\item[BiLSTM] Bidirectional Long Short-Term Memory (Memória de Longo e Curto Prazo Bidirecional)
\item[CDSS] Clinical Decision Support Systems (Sistemas de Suporte à Decisão Clínica)
\item[CHB-MIT] Children's Hospital Boston - Massachusetts Institute of Technology
\item[CNN] Convolutional Neural Network (Rede Neural Convolucional)
\item[CPF] Cadastro de Pessoas Físicas
\item[CRM] Conselho Regional de Medicina
\item[EDF] European Data Format
\item[EEG] Eletroencefalograma
\item[eMAG] Modelo de Acessibilidade em Governo Eletrônico
\item[FROC] Free-Response Receiver Operating Characteristic
\item[HTTP] Hypertext Transfer Protocol
\item[IA] Inteligência Artificial
\item[IFG] Instituto Federal de Educação, Ciência e Tecnologia de Goiás
\item[I/O] Input/Output (Entrada/Saída)
\item[IoU] Intersection over Union (Interseção sobre União)
\item[JSON] JavaScript Object Notation
\item[JWT] JSON Web Token
\item[LGPD] Lei Geral de Proteção de Dados Pessoais
\item[LSTM] Long Short-Term Memory (Memória de Longo e Curto Prazo)
\item[MER] Modelo Entidade-Relacionamento
\item[MLP] Multi-Layer Perceptron (Perceptron Multicamadas)
\item[P\&D] Pesquisa e Desenvolvimento
\item[PIBITI] Programa Institucional de Bolsas de Iniciação em Desenvolvimento Tecnológico e Inovação
\item[RDC] Resolução da Diretoria Colegiada
\item[REST] Representational State Transfer
\item[RF] Requisito Funcional
\item[RNF] Requisito Não Funcional
\item[SADC] Sistema de Apoio à Decisão Clínica
\item[SHAP] Shapley Additive exPlanations
\item[SOA] Service-Oriented Architecture (Arquitetura Orientada a Serviços)
\item[SPA] Single Page Application
\item[UI] User Interface (Interface de Usuário)
\item[WCAG] Web Content Accessibility Guidelines
\item[WHO] World Health Organization (Organização Mundial da Saúde)
\item[XAI] Explainable Artificial Intelligence (Inteligência Artificial Explicável)
\end{siglas}

\pdfbookmark[0]{\contentsname}{toc}
\tableofcontents*
\cleardoublepage
